\begin{lemma}
Let $f,g\colon\mathbb{C}\to\mathbb{C}$ be real-differentiable at
$\xi\in\mathbb{C}$.  Then
\[
  \frac{\partial(fg)}{\partial\xi}(\xi)
  =\frac{\partial f}{\partial\xi}(\xi)\,g(\xi)+f(\xi)\,\frac{\partial g}{\partial\xi}(\xi),
\]
\[
  \frac{\partial(fg)}{\partial\bar\xi}(\xi)
  =\frac{\partial f}{\partial\bar\xi}(\xi)\,g(\xi)+f(\xi)\,\frac{\partial g}{\partial\bar\xi}(\xi),
\]
where $fg$ denotes the pointwise product $z\mapsto f(z)g(z)$,
$\partial/\partial\xi$ is as in \noderef{DZ} and $\partial/\partial\bar\xi$ is
as in \noderef{DZbar}.
\end{lemma}
\begin{proof}
For a map $h\colon\mathbb{C}\to\mathbb{C}$ real-differentiable at $\xi$ write
$Dh(\xi)$ for its real Fr\'echet derivative at $\xi$, a real-linear map
$\mathbb{C}\to\mathbb{C}$.  By \noderef{DZ} and \noderef{DZbar},
\begin{equation*}
  \frac{\partial h}{\partial\xi}(\xi)=\tfrac12\bigl(Dh(\xi)[1]-i\,Dh(\xi)[i]\bigr),
  \qquad
  \frac{\partial h}{\partial\bar\xi}(\xi)=\tfrac12\bigl(Dh(\xi)[1]+i\,Dh(\xi)[i]\bigr).
\end{equation*}

Multiplication $\mathbb{C}\times\mathbb{C}\to\mathbb{C}$ is a continuous
$\mathbb{R}$-bilinear map, so since $f$ and $g$ are real-differentiable at
$\xi$ the pointwise product $fg$ is real-differentiable at $\xi$ and its
derivative is given by the bilinear product rule
\begin{equation}\tag{$*$}
  D(fg)(\xi)[v]=\bigl(Df(\xi)[v]\bigr)g(\xi)+f(\xi)\bigl(Dg(\xi)[v]\bigr)
  \qquad\text{for every }v\in\mathbb{C}.
\end{equation}

Insert $(*)$ with $v=1$ and with $v=i$ into the formula for
$\partial/\partial\xi$:
\begin{align*}
  \frac{\partial(fg)}{\partial\xi}(\xi)
  &=\tfrac12\Bigl(D(fg)(\xi)[1]-i\,D(fg)(\xi)[i]\Bigr)\\
  &=\tfrac12\Bigl(\bigl(Df(\xi)[1]\bigr)g(\xi)+f(\xi)\bigl(Dg(\xi)[1]\bigr)
      -i\bigl(Df(\xi)[i]\bigr)g(\xi)-i\,f(\xi)\bigl(Dg(\xi)[i]\bigr)\Bigr)\\
  &=\tfrac12\bigl(Df(\xi)[1]-i\,Df(\xi)[i]\bigr)\,g(\xi)
    +f(\xi)\,\tfrac12\bigl(Dg(\xi)[1]-i\,Dg(\xi)[i]\bigr)\\
  &=\frac{\partial f}{\partial\xi}(\xi)\,g(\xi)+f(\xi)\,\frac{\partial g}{\partial\xi}(\xi),
\end{align*}
where the third equality only regroups the four terms, using that $g(\xi)$ and
$f(\xi)$ are fixed complex numbers and that multiplication in $\mathbb{C}$ is
commutative and distributive.

The second identity is the same computation with $+i$ in place of $-i$
throughout:
\begin{align*}
  \frac{\partial(fg)}{\partial\bar\xi}(\xi)
  &=\tfrac12\Bigl(D(fg)(\xi)[1]+i\,D(fg)(\xi)[i]\Bigr)\\
  &=\tfrac12\bigl(Df(\xi)[1]+i\,Df(\xi)[i]\bigr)\,g(\xi)
    +f(\xi)\,\tfrac12\bigl(Dg(\xi)[1]+i\,Dg(\xi)[i]\bigr)\\
  &=\frac{\partial f}{\partial\bar\xi}(\xi)\,g(\xi)+f(\xi)\,\frac{\partial g}{\partial\bar\xi}(\xi).
\end{align*}
\end{proof}
